%% ma: L10-B01-0027
%% lop: 10
%% chuong: 1
%% bai: 1
%% dang: 10.1.6
%% loai: TN4
%% muc: TH
%% dap_an: "D"
%% nguon: "Ảnh giáo viên gửi 10/10/2026 (ThS. Nguyễn Hoàng Việt, Chương 1 Mệnh đề), B-Đề 2, Phần 1 Câu 8"
%% kien_thuc: "Bài 1 (mệnh đề kéo theo, mệnh đề tương đương)"
%% trang_thai: cho-duyet
%% ly_do: "Dạng toán mới đề xuất 10.1.6: mệnh đề kéo theo, tương đương: xét tính đúng sai"
%% ghi_chu: "Đáp án dự kiến D. Xét chặt chẽ, chiều đảo của B và C chỉ đúng khi H nằm giữa B và C (tam giác tù vẫn có thể thỏa BA^2=BH.BC); giáo viên cân nhắc thêm điều kiện này vào đề"
\begin{ex}
Cho tam giác $ABC$ với $H$ là chân đường cao từ $A$. Mệnh đề nào sau đây sai?
\choice{``$ABC$ là tam giác vuông ở $A\Leftrightarrow\dfrac1{AH^2}=\dfrac1{AB^2}+\dfrac1{AC^2}$''}{``$ABC$ là tam giác vuông ở $A\Leftrightarrow BA^2=BH\cdot BC$''}{``$ABC$ là tam giác vuông ở $A\Leftrightarrow HA^2=HB\cdot HC$''}{\True ``$ABC$ là tam giác vuông ở $A\Leftrightarrow BA^2=BC^2+AC^2$''}
\loigiai{Các hệ thức ở A, B, C là các hệ thức lượng trong tam giác vuông tại $A$ (với đường cao $AH$). D sai: hệ thức $BA^2=BC^2+AC^2$ là định lí Pythagore cho tam giác vuông tại $C$ (cạnh huyền $BA$), không tương đương với $\widehat A=90^\circ$ (khi $\widehat A=90^\circ$ thì $BC^2=BA^2+AC^2$).}
\end{ex}
